\documentclass[a4paper]{scrartcl}
\usepackage[utf8]{inputenc}
\usepackage{amsmath,amssymb,amsthm}
\usepackage{enumitem,setspace}
\usepackage[usenames,dvipsnames,table]{xcolor}
\usepackage{graphicx}
\DeclareGraphicsExtensions{.pdf,.png,.eps,.jpg}
\usepackage{algorithm2e}

% --- math operators ---
\usepackage{mathtools}
\DeclarePairedDelimiter\set{\{}{\}} % use $\set{1, 2, 3}$
\DeclarePairedDelimiter\abs{\lvert}{\rvert}
\DeclarePairedDelimiter\norm{\lVert}{\rVert}
\DeclarePairedDelimiter\ceils{\lceil}{\rceil}
\DeclarePairedDelimiter\floor{\lfloor}{\rfloor}
\def\then{\ensuremath{\Rightarrow}} % =>
\def\iff{\ensuremath{\Leftrightarrow}} % if and only if <=>
\def\to{\ensuremath{\rightarrow}} % ->
\def\ot{\ensuremath{\leftarrow}} % ->
\def\Oh{\ensuremath{\mathcal{O}}} % big O like $\Oh(n)$
% ---
\DeclareMathOperator{\preorder}{preorder}
\DeclareMathOperator{\postorder}{postorder}
\DeclareMathOperator{\depth}{depth}

% --- FILL OUT THESE 3 FIELDS ---
\def\sheetNumber{1}
\def\names{NAMEN} 
\def\sumPoints{10} 

% --- ### BEGIN DOCUMENT ### ---
\begin{document} 

\begin{doublespace} % header
\noindent \textbf{\LARGE{Übungsblatt \sheetNumber}}
\hfill  {\large Punkte: $\boxed{\qquad  /\; \sumPoints}$}\\
{\Large \textbf{Fortgeschrittene Algorithmen (Wintersemester 2022/23)}\\
Bearbeitet von: \textbf{\names}\\
\today}
\end{doublespace}


\section*{Aufgabe 1}

a) Behauptung. Sei \(G=(V,E,w)\) ein ungerichteter, gewichteter Graph mit positiven Kantengewichten und sei
\(e=\{u,v\}\in E\) die eindeutig leichte Kante, d.h. \(w(e)<w(e')\) für alle \(e'\in E\setminus\{e\}\).
Dann ist \(e\) in jedem minimalen Spannbaum (MST) von \(G\) enthalten.\\

Widerspruchsbeweis: Angenommen, es existiere ein minimaler Spannbaum \(T\) von \(G\) mit \(e\notin T\).
Durch Hinzufügen von \(e\) zu \(T\) entsteht ein Zyklus \(C\subseteq T\cup\{e\}\), der natürlich \(e\) enthält.
Wegen der Eindeutigkeit von \(e\) gilt für jede andere Kante \(f\in C\setminus\{e\}\) die Ungleichung \(w(f)>w(e)\).
Entfernt man nun eine solche Kante \(f\) aus \(C\), so erhält man wieder einen Spannbaum
\(T' = T + e - f\). Für die Gewichte folgt
\[
w(T') = w(T) + w(e) - w(f) < w(T),
\]
da \(w(e)<w(f)\). Damit ist \(T'\) ein Spannbaum mit strikt geringerem Gesamtgewicht als \(T\), was der Annahme widerspricht, dass \(T\) minimal sei.

Folglich kann die Annahme nicht gelten und \(e\) muss in jedem minimalen Spannbaum von \(G\) enthalten sein.

% -- code for figures
% \begin{figure}[htb]
% \centering
% \includegraphics{filename} % use option [width=\linewidth] to change size to linewidth (or 0.6\linewidth for example)
% \label{fig:name}
% \end{figure}


\end{document}
% --- ### END DOCUMENT ### ---